\section{Grundlagen}

\subsection{Gleichungen zur Modellbildung}

\begin{tcolorbox}
\boxtitle{Kräftegleichungen (Translation)}
Federkraft: $F_F = k_F \cdot x$\\[2pt]
Dämpferkraft: $F_D = k_D \cdot v = k_D \cdot \dot{x}$\\[2pt]
Trägheitskraft: $F_{Tr} = m \cdot a = m \cdot \ddot{x}$\\[2pt]
Erdanziehungskraft: $F_G = m \cdot g$
\end{tcolorbox}

\begin{tcolorbox}
\boxtitle{Momentengleichungen (Rotation)}
Widerstandsmoment: $M_W = k_D \cdot \omega$\\[2pt]
Trägheitsmoment: $M_{Tr} = J \cdot \dot{\omega} = J \cdot \ddot{\varphi}$
\end{tcolorbox}

\begin{tcolorbox}
\boxtitle{Elektrischer Grundkreis (RLC)}
\[
u(t) = L \cdot \frac{di}{dt} + R \cdot i + \frac{1}{C} \int i \, dt
\]
\end{tcolorbox}

\begin{tcolorbox}
\boxtitle{Rotation in Flüssigkeit}
\[
M = M_{\text{Träg}} + M_{\text{Brems}} = J \cdot \dot{\omega} + k_{\text{Flüssigkeit}} \cdot \omega
\]
\end{tcolorbox}

\subsection{Mathematische Grundlagen}
\begin{tcolorbox}
\boxtitle{Kleine Winkel}
Für kleine Winkel $\alpha$ (im Bogenmaß) gilt:\\[2pt]
$\sin(\alpha) \approx \alpha, \quad \tan(\alpha) \approx \alpha, \quad \cos(\alpha) \approx 1 - \frac{\alpha^2}{2} \approx 1$
\end{tcolorbox}

\begin{tcolorbox}
\boxtitle{Grenzwertsätze der Laplace-Transformation}
\textbf{Anfangswertsatz:}\\[2pt]
\[
x(t \to +0) = \lim_{s \to \infty} s \cdot X(s)
\]
\textbf{Endwertsatz:}\\[2pt]
\[
x(t \to \infty) = \lim_{s \to 0} s \cdot X(s)
\]
\textbf{Gültigkeitsbedingung:} Der Endwertsatz gilt \textbf{nur}, wenn höchstens ein Pol von $s \cdot X(s)$ im Ursprung liegt und alle anderen Pole in der offenen linken $s$-Halbebene liegen ($\text{Re}\{s_P\} < 0$).
\end{tcolorbox}

\subsection{Linearisierung um den Arbeitspunkt}
\begin{tcolorbox}
Nichtlineare Kennlinie $x_a = f(x_{e1}, \dots, x_{en})$ um den stationären Arbeitspunkt $AP = (x_{e1,AP}, \dots, x_{en,AP})$:
\begin{align*}
x_a(t) &= x_{a,AP} + \Delta x_a(t) \\
&\approx f(x_{e1,AP}, \dots, x_{en,AP}) + \sum_{m=1}^n \left. \frac{\partial f}{\partial x_{em}} \right|_{AP} \cdot \Delta x_{em}(t)
\end{align*}
Lineare Kleinsignalbeziehung für die Abweichungen:
\[
\Delta x_a(t) \approx \sum_{m=1}^n K_m \cdot \Delta x_{em}(t) \quad \text{mit } K_m = \left. \frac{\partial f}{\partial x_{em}} \right|_{AP}
\]
\end{tcolorbox}

\subsection{Eigenschaften von LTI-Systemen}
\begin{tcolorbox}
\begin{enumerate}
  \item \textbf{BIBO-Stabilität} (Bounded Input / Bounded Output):
  \[
  |y(t)| < \infty \quad \text{für } |x(t)| < \infty
  \]
  \item \textbf{Linearität} (Superposition und Homogenität):
  \[
  f\left\{ \sum_{n=1}^N a_n x_n(t) \right\} = \sum_{n=1}^N a_n f\{x_n(t)\}
  \]
  \item \textbf{Zeitinvarianz}:
  \[
  f\{x(t - t_0)\} = y(t - t_0)
  \]
  \item \textbf{Kausalität}: Das System reagiert nicht vor Eintreffen des Eingangssignals:
  \[
  x(t) = 0 \text{ für } t < 0 \implies y(t) = 0 \text{ für } t < 0
  \]
\end{enumerate}
\end{tcolorbox}

\subsection{Systemantworten}
\begin{tcolorbox}
\boxtitle{Sprungantwort ($x(t) = \sigma(t)$)}
\[
g(t) = \mathcal{L}^{-1}\{G(s)\} = \mathcal{L}^{-1}\left\{ \frac{F(s)}{s} \right\} = \int_0^t h(\tau) \, d\tau
\]
\end{tcolorbox}

\begin{tcolorbox}
\boxtitle{Impulsantwort ($x(t) = \delta(t)$)}
\[
h(t) = \int_{-\infty}^\infty x(\tau) \delta(t - \tau) \, d\tau = \mathcal{L}^{-1}\{F(s)\} = \frac{d}{dt} g(t)
\]
\end{tcolorbox}

\subsection{Abtasttheorem}
\begin{tcolorbox}
Durch die periodische Abtastung eines Signals $f(t)$ mit der Periode $T_A$ (bzw. $\omega_A = 2\pi/T_A$) wird das kontinuierliche Spektrum $F(\omega)$ unendlich oft um die Vielfachen der Abtastkreisfrequenz reproduziert:
\[
F_A(\omega) = \frac{1}{T_A} \sum_{n=-\infty}^\infty F(\omega - n \omega_A)
\]
\textbf{Shannon/Nyquist-Bedingung:}\\[2pt]
Zur fehlerfreien (aliasingfreien) Rekonstruktion muss die Abtastfrequenz mindestens doppelt so groß wie die höchste Grenzfrequenz $\omega_g$ des abzutastenden Signals sein:
\[
\omega_A \ge 2 \omega_g \quad \text{bzw.} \quad f_A \ge 2 f_g
\]
\end{tcolorbox}
